\title{QLoRA: Efficient Finetuning of Quantized LLMs}

\begin{abstract}
We introduce \textsc{QLoRA}, a memory-efficient finetuning method that makes it possible to adapt a 65B parameter language model on a single 48GB GPU, all while maintaining the finetuning effectiveness seen with standard 16-bit precision. In \textsc{QLoRA}, gradients are backpropagated through a fixed, 4-bit quantized pretrained model and into Low Rank Adapters~(LoRA). Our premier model family, \textbf{Guanaco}, surpasses previously published open-access models on the Vicuna benchmark, attaining 99.3\% of ChatGPT’s benchmarked performance after just 24 hours of finetuning on one GPU. Key innovations in \textsc{QLoRA} that facilitate memory efficiency without performance degradation include: (a) the introduction of 4-bit NormalFloat (NF4), a data format theoretically optimized for representing normally distributed parameters, (b) Double Quantization, which further decreases memory requirements by quantizing the quantization parameters themselves, and (c) Paged Optimizers, which address and alleviate transient memory demands. Utilizing \textsc{QLoRA}, we conduct finetuning experiments on over 1,000 models and provide a thorough evaluation across 8 instruction datasets, various architectures (such as LLaMA, T5), and large model sizes (e.g., 33B and 65B parameters), which are otherwise impractical for standard finetuning approaches. The findings reveal that \textsc{QLoRA}-enabled finetuning on compact, high-quality datasets achieves or exceeds state-of-the-art results, even with smaller models than earlier leading systems. We present an in-depth comparison of chatbot abilities based on both human judgments and GPT-4-based assessments, and demonstrate that GPT-4 can serve as a cost-effective and reliable substitute for human evaluation. Our study also highlights the limitations and unreliability of current chatbot benchmarks for assessing model quality. A targeted analysis illustrates specific cases where \textbf{Guanaco} does not match ChatGPT's performance. We make available our models and the associated codebase, including CUDA implementations for 4-bit training.\footnote{\url{https://github.com/artidoro/qlora} and \url{https://github.com/TimDettmers/bitsandbytes}}
\end{abstract}

\section{Introduction}

Refining large language models (LLMs) through finetuning is widely recognized as an effective strategy for enhancing their capabilities \citep{min2021metaicl, wei2021finetuned, ouyang2022training, wang2022super, wang2022self, liu2022few}, as well as for instilling preferred behaviors or suppressing unwanted ones \citep{ouyang2022training,askell2021general,bai2022training}. Nevertheless, the resource demands for finetuning extremely large models are often prohibitive; for example, conventional 16-bit finetuning of the LLaMA model with 65B parameters~\citep{touvron2023llama} necessitates in excess of 780 GB of GPU memory. While recent advances in quantization have enabled a reduction in the inference-time memory consumption for LLMs \citep{dettmers2022llm, dettmers2022case, frantar2022gptq, xiao2022smoothquant}, these methods are not applicable during the training process, as their stability breaks down \citep{wortsman2023stable}.

In this work, we establish for the first time that finetuning a model quantized to 4-bits can be achieved without incurring any loss in model performance. Our approach, \textsc{QLoRA}, leverages an innovative high-precision quantization scheme to transform a pretrained network to 4-bit precision, followed by the integration of a compact set of trainable Low-rank Adapter weights \citep{hu2021lora}. These adapters are updated through gradient backpropagation across the quantized parameters.

Utilizing \textsc{QLoRA}, the memory required to finetune a 65B parameter network drops from more than 780GB to under 48GB of GPU memory, with no compromise in speed or accuracy relative to standard 16-bit finetuning. This dramatically improves access to LLM finetuning, enabling even the largest open-source models to be finetuned using just a single GPU. With \textsc{QLoRA}, we introduce the \textbf{Guanaco} series of models: the second-best variant achieves 97.8\% of ChatGPT’s performance on the Vicuna~\citep{vicuna2023} benchmark and can be finetuned within 12 hours on a standard consumer GPU; using a single high-end GPU over 24 hours, our largest model attains 99.3\%, thus nearly matching ChatGPT’s performance. The smallest \textbf{Guanaco} model (7B parameters) operates with only 5 GB of memory while surpassing the 26 GB Alpaca model by over 20 percentage points on the Vicuna benchmark (see Table~\ref{tbl:vicuna_eval}).

\textsc{QLoRA} incorporates several key advancements aimed at minimizing memory consumption while maintaining model performance: (1) \textbf{4-bit NormalFloat}, a quantization data type that is theoretically optimal for data with a normal distribution and empirically outperforms both 4-bit Integers and 4-bit Floats. (2) \textbf{Double Quantization}, which applies quantization to the quantization constants themselves, yielding an average reduction of approximately 0.37 bits per parameter (around 3 GB for models with 65B parameters). (3) \textbf{Paged Optimizers}, leveraging NVIDIA unified memory to prevent the abrupt memory increases associated with gradient checkpointing, particularly for mini-batches involving lengthy sequences. By integrating these techniques, we enhance the LoRA method by including adapters at every layer of the network, thereby significantly reducing the accuracy compromises observed in previous research.

The resource efficiency of \textsc{QLoRA} enables comprehensive exploration of instruction fine-tuning and chatbot capabilities across model sizes that would be infeasible with traditional fine-tuning methods due to memory constraints. Consequently, we fine-tune over 1,000 models spanning various instruction tuning datasets, architectures, and parameter counts ranging from 80M to 65B. Beyond demonstrating that \textsc{QLoRA} can match 16-bit performance (\S\ref{sec:comp_to_std}) and facilitate the training of the cutting-edge chatbot \textbf{Guanaco} (\S\ref{sec:pushing}), we further investigate patterns emerging from the trained models. Our analysis reveals two main insights. First, the quality of training data plays a substantially larger role than the volume of data—for instance, a 9k-sample dataset (OASST1) surpassed a much larger 450k-sample dataset (FLAN v2, subsampled) in chatbot evaluations, even when both are used for instruction following. Second, high scores on the Massive Multitask Language Understanding (MMLU) benchmark do not necessarily translate into strong performance on the Vicuna chatbot benchmark, and vice versa. This suggests that, for a particular task, the appropriateness of the dataset is more crucial than its scale.

Additionally, we conduct a thorough assessment of chatbot effectiveness utilizing both human annotators and GPT-4 for evaluation purposes. Our methodology involves tournament-style benchmarking, where models directly compete by generating responses to the same prompts, and either GPT-4 or human evaluators determine which response is superior in each match. We then aggregate tournament outcomes into Elo scores~\citep{elo1967proposed, elo1978rating}, which provide a systematic ranking of model performance. Our findings indicate that the rankings assigned by GPT-4 and by human judges generally correspond, though we also observe notable discrepancies in some cases. Thus, we caution that while model-based evaluation presents a scalable and cost-efficient substitute for human annotation, it is accompanied by inherent sources of unreliability.

To complement our quantitative benchmarks, we also present a qualitative examination focused on the \textbf{Guanaco} models. This analysis showcases both the strengths and weaknesses of the models, particularly highlighting areas not sufficiently captured by numerical evaluation metrics.

In support of ongoing research, we make available all generated responses from models, alongside corresponding human and GPT-4 annotations. Furthermore, we release our implementation codebase and CUDA kernels, integrating our approach within the Hugging Face transformers library \citep{wolf2019huggingface} to maximize accessibility. Our release includes a suite of adapters for models with 7/13/33/65B parameters, each trained on eight different instruction-following datasets, resulting in a total of 32 open-source, finetuned models.

\section{Background}
\paragraph{Block-wise k-bit Quantization} Quantization refers to converting data representations with higher precision into lower-precision formats. This commonly involves reducing the bit width, such as transforming 32-bit floating-point numbers into 8-bit integers. To utilize the complete dynamic range of the target lower-bit representation, the data is usually normalized with respect to the maximum absolute value found in the original tensor, thus rescaling it before quantization. For instance, to quantize a tensor of 32-bit floating-point values (FP32) into 8-bit integers (Int8) constrained to the range $[-127, 127]$:

\begin{equation}
    \mathbf{X}^{ \text{Int8}} = \text{round}\left( \frac{127}{{\text{absmax}}(\mathbf{X}^{{\text{FP32}}})}\mathbf{X}^{{\text{FP32}}} \right) = \text{round}(c^{\text{FP32}}\cdot \mathbf{X}^{{\text{FP32}}}),
\end{equation}

where $c$ denotes the {\em quantization constant} or {\em quantization scale}. To revert to the original format, dequantization is applied as follows:

\begin{equation}
    \text{dequant}(c^{\text{FP32}}, \mathbf{X}^{\text{Int8}}) = \frac{\mathbf{X}^{\text{Int8}}}{c^{\text{FP32}}} = \mathbf{X}^{\text{FP32}}
\end{equation}

This method faces challenges when the input tensor includes a value with an unusually large magnitude (i.e., an outlier). In such cases, many quantization bins (bit configurations) are poorly utilized, with few or no input values mapped to certain bins. To alleviate the problem caused by outliers, a widely used solution involves partitioning the input tensor into multiple blocks, each quantized independently and assigned a distinct quantization constant $c$. More precisely, the input tensor $\mathbf{X}\in \mathbb{R}^{b\times h}$ is flattened and divided into $n$ consecutive blocks of size $B$, so that the tensor is split into ${n = ({b\times h} )/{B} }$ segments. Each segment is then quantized separately using Equation~1, resulting in a quantized tensor accompanied by $n$ individual quantization constants $c_i$.

\paragraph{Low-rank Adapters} The Low-rank Adapter (LoRA) finetuning technique \citep{hu2021lora} addresses memory constraints by introducing a small set of trainable adapter parameters, while the main model weights are kept static. During training via stochastic gradient descent, the gradient signals propagate through the fixed pretrained model parameters to update only the adapter parameters, with the goal of minimizing the loss. LoRA implements this by supplementing a linear projection with an additional low-rank projection. Given a linear mapping ${\mathbf{X} \mathbf{W} = \mathbf{Y}}$ where $\mathbf{X}\in  \mathbb{R}^{b\times h}$ and $\mathbf{W}\in  \mathbb{R}^{h\times o}$, LoRA produces:

\begin{equation}
    \mathbf{Y} = \mathbf{X}\mathbf{W} + s\mathbf{X}\mathbf{L}_1\mathbf{L}_2,
\end{equation}

with $\mathbf{L}_1\in  \mathbb{R}^{h\times r}$, $\mathbf{L}_2\in  \mathbb{R}^{r\times o}$, and a scaling factor $s$.

\paragraph{Memory Requirement of Parameter-Efficient Finetuning} A key consideration is the memory consumption of LoRA during training, both in relation to how many adapters are employed and their respective sizes. Due to the extremely lightweight nature of LoRA's memory demands, it is feasible to scale up the number of adapters to boost performance without causing a significant rise in overall memory utilization. Although LoRA is categorized as a Parameter Efficient Finetuning (PEFT) strategy, the bulk of memory usage in LLM finetuning originates from storing activation gradients rather than from the LoRA parameters themselves. For instance, when fine-tuning a 7B LLaMA model on FLAN v2 with a batch size of 1, and using LoRA weights set to 0.2\% of the base model’s parameters—as commonly recommended in the literature~\citep{hu2021lora,liu2022few}—the memory footprint for the LoRA input gradients reaches 567 MB, in contrast to just 26 MB for the LoRA parameters themselves. Applying gradient checkpointing~\citep{chen2016training} further brings down the memory required for input gradients to an average of 18 MB per sequence, but this still surpasses the total size of the LoRA parameters. Meanwhile, the 4-bit quantized model’s base weights alone require 5,048 MB of memory. This demonstrates not only the significance of using gradient checkpointing but also that further reducing LoRA parameter sizes results in only minimal savings in total memory usage. Therefore, increasing the number of adapters has little impact on the memory budget during training (a comprehensive breakdown is available in Appendix~\ref{app:memory}). As will be elaborated upon later, this characteristic is particularly important for matching the performance of full 16-bit precision.

\section{\textsc{QLoRA} Finetuning}

\textsc{QLoRA} enables high-fidelity finetuning at 4-bit precision through two primary innovations: 4-bit NormalFloat (NF4) quantization and Double Quantization. We further present Paged Optimizers, a strategy aimed at circumventing memory surges during gradient checkpointing, which have historically led to out-of-memory issues when finetuning large-scale models on a single machine.

Within the \textsc{QLoRA} framework, a single low-precision format—commonly 4-bit—is employed for storage, alongside a higher-precision data type (typically BFloat16) for computation. Operationally, each time a \textsc{QLoRA} parameter tensor is accessed, it undergoes dequantization to BFloat16 prior to performing 16-bit matrix multiplications.

Below, we elaborate on the elements of \textsc{QLoRA} before proceeding to its formal specification.

\paragraph{4-bit NormalFloat Quantization} The NormalFloat (NF) format extends the concept of Quantile Quantization \citep{dettmers2022optimizers}, an information-theoretically optimal representation where each bin is assigned the same number of entries from the original tensor. Quantile quantization operates by estimating input tensor quantiles based on the empirical cumulative distribution function.

A significant drawback of quantile quantization lies in the high computational expense required for accurate quantile estimation. Thus, fast but approximate methods such as SRAM quantiles~\citep{dettmers2022optimizers} are commonly used. However, these approximations can introduce substantial quantization errors, particularly for outlier values—which frequently have the greatest significance.

If input tensors are drawn from a distribution that remains unchanged except for a quantization constant, quantile estimation becomes much more efficient. In such circumstances, because the quantiles are consistent across tensors, precise quantile computation is both possible and computationally manageable.

Pretrained neural network weights are typically distributed according to a normal distribution centered at zero with a standard deviation $\sigma$ (refer to Appendix~\ref{app:normality}). To map these weights onto a consistent, predetermined distribution, we adjust $\sigma$ so that the resulting distribution is confined entirely within the data type's specified range. In this context, we define the data type's interval as $[-1, 1]$. Therefore, both the quantization levels for the data type and the original neural network weights must be standardized to this $[-1, 1]$ interval.

For zero-mean normal distributions of arbitrary standard deviation $\sigma$ that are mapped into the range $[-1, 1]$, the information-theoretically optimal quantization scheme is established through the following procedure: (1) compute the $2^k+1$ quantiles of an idealized $N(0, 1)$ distribution to define a $k$-bit quantile-based quantization format appropriate for normal distributions; (2) rescale the resulting data type's values so that they span the interval $[-1, 1]$; and (3) quantize the neural network weight tensor by scaling it into $[-1, 1]$ via normalization by its absolute maximum.

With the support of both the data type and the weights thus aligned, standard quantization methods may be directly applied. The third step described here effectively rescales the weight tensor's standard deviation so it matches that of the $k$-bit quantized type. Formally, the $2^k$ quantized values $q_i$ are determined by

\begin{equation}
q_i = \frac{1}{2}\left( Q_X\left(\frac{i}{2^k+1}\right) + Q_X\left(\frac{i+1}{2^k+1}\right)\right),
\end{equation}

where $Q_X(\cdot)$ denotes the quantile function of the standard normal distribution $N(0,1)$. One limitation of symmetric $k$-bit quantization is the lack of an exact representation for zero, which is essential for accurately encoding padding and other elements with zero value. To guarantee that zero is mapped to a discrete quantization level and to fully utilize all $2^k$ distinct values in a $k$-bit data type, we devise an asymmetric quantization scheme. Specifically, we determine quantiles $q_i$ corresponding to $2^{k-1}$ levels for the negative side and $2^{k-1}+1$ levels for the positive side, after which we combine the two sets of $q_i$, discarding one of the duplicate zeros that appears in both. The resulting data type, which maintains a uniform expected distribution of elements within each bin, is referred to as the \emph{k-bit NormalFloat} (NFk). This construction is information-theoretically optimal for data drawn from zero-centered normal distributions. The precise definition and quantized values for this type can be found in Appendix~\ref{app:nf4}.

\paragraph{Double Quantization{}} 
We propose the method of \emph{Double Quantization} (DQ), in which the quantization constants themselves are further quantized to achieve additional memory efficiency. Although accurate 4-bit quantization~\citep{dettmers2022case} typically requires a small blocksize, this comes with a non-trivial memory cost: for example, employing 32-bit quantization constants and a blocksize of 64 for $\mathbf{W}$ incurs an average storage overhead of $32/64 = 0.5$ bits per parameter. Double Quantization addresses this issue by compressing the quantization constants themselves.

To implement this, we perform a secondary quantization step on the original quantization constants $c_2^{\text{FP32}}$. These constants are quantized to $c_2^{\text{FP8}}$ using 8-bit floating point representation, along with an associated set of higher-level quantization constants $c_1^{\text{FP32}}$. The secondary quantization employs a blocksize of 256, leveraging the observation from \citet{dettmers2022case} that 8-bit quantization does not harm performance. Since the $c_2^{\text{FP32}}$ are inherently positive, we mean-center them prior to quantization, enabling symmetric quantization schemes to be used. This process reduces the average memory cost per parameter from $0.5$ bits (i.e., $32/64$) to $8/64 + 32/(64\cdot256) = 0.127$ bits for a blocksize of 64, resulting in a savings of $0.373$ bits per parameter.

\paragraph{Paged Optimizers} leverage the NVIDIA unified memory~\footnote{\tiny \url{https://docs.nvidia.com/cuda/cuda-c-programming-guide}} capability, which facilitates seamless, on-demand transfers of memory pages between the CPU and GPU. This mechanism is particularly advantageous when GPU memory is temporarily exhausted, enabling uninterrupted GPU computation by automating data migration. Analogous to conventional memory paging between main memory and disk, unified memory handles data movement transparently. In our approach, we allocate the optimizer state variables in paged memory, allowing these states to be automatically migrated to CPU memory in response to GPU memory pressure, and then recalled to GPU memory during the optimizer’s update step as required.

\paragraph{\textsc{QLoRA}.} With the preceding components in place, we define \textsc{QLoRA} for a quantized base model’s linear layer with a LoRA adapter as:

\begin{equation}
    \mathbf{Y}^{\text{BF16}} =  \mathbf{X}^{\text{BF16}} \text{doubleDequant}(c_1^{\text{FP32}}, c_2^{\text{k-bit}}, \mathbf{W}^{\text{NF4}}) + \mathbf{X}^{\text{BF16}}\mathbf{L}^{\text{BF16}}_1\mathbf{L}^{\text{BF16}}_2,
\end{equation}

where the doubleDequant$(\cdot)$ operator is given by:

\begin{equation}
    \text{doubleDequant}(c_1^{\text{FP32}}, c_2^{\text{k-bit}}, \mathbf{W}^{\text{k-bit}}) = \text{dequant}(\text{dequant}(c_1^{\text{FP32}}, c_2^{\text{k-bit}}), \mathbf{W}^{\text{4bit}}) = \mathbf{W}^{\text{BF16}},
\end{equation}

Here, we utilize NF4 quantization for $\mathbf{W}$ and represent $c_2$ in FP8 format.

To achieve a higher degree of quantization accuracy for $\mathbf{W}$, we set its blocksize to 64; for $c_2$, a blocksize of 256 is chosen to minimize memory usage.

In the optimization step, only gradients of the error with respect to the adapter parameters, $\frac{\partial E}{\partial \mathbf{L}_i}$, are necessary—gradients with respect to the quantized weights $\frac{\partial E}{\partial \mathbf{W}}$ are not computed. However, evaluating $\frac{\partial E}{\partial \mathbf{L}_i}$ still requires computing $\frac{\partial \mathbf{X}}{\partial \mathbf{W}}$, which relies on Eq.~(5). This involves dequantizing $\mathbf{W}^{\text{NF4}}$ to the computational type $\mathbf{W}^{\text{BF16}}$, ensuring the derivative $\frac{\partial \mathbf{X}}{\partial \mathbf{W}}$ is obtained in BFloat16 precision.

In summary, \textsc{QLoRA} employs a single data type for storage, typically 4-bit NormalFloat, alongside a computation data type, which is 16-bit BrainFloat. During both the forward and backward passes, the stored values are dequantized into the computation data type; however, only the LoRA parameters, maintained in 16-bit BrainFloat, have their weight gradients computed.

\section{QLoRA vs. Standard Finetuning}
\label{sec:comp_to_std}

We have outlined the mechanics of QLoRA and highlighted its substantial memory savings in the context of finetuning large models. An essential consideration, however, is whether QLoRA’s performance matches that of conventional full-model finetuning. Moreover, we are interested in dissecting specific QLoRA components, such as the effectiveness of NormalFloat4 relative to standard Float4. The experimental results reported in the next sections address these concerns.

\paragraph{Experimental setup.} We examine three types of model architectures—encoder, encoder-decoder, and decoder-only—and benchmark QLoRA against both 16-bit adapter-based finetuning and complete finetuning for models of up to 3B parameters. Our evaluation spans several benchmarks: GLUE \citep{wang2018glue} with RoBERTa-large \citep{liu2019roberta}, Super-NaturalInstructions (TKInstruct) \citep{wang2022super} with T5 \citep{t5}, and 5-shot MMLU \citep{hendrycksmeasuring} subsequent to finetuning LLaMA on Flan v2 \citep{longpre2023flan} and Alpaca \citep{alpaca}. To investigate NF4’s benefits over alternative 4-bit formats, we replicate the methodology of \citet{dettmers2022case}, measuring both post-quantization zero-shot accuracy and perplexity across several model families (OPT~\citep{zhang2022opt}, LLaMA~\citep{touvron2023llama}, BLOOM~\citep{scao2022bloom}, Pythia~\citep{biderman2023pythia}) spanning model sizes from 125m to 13B.

Each result section elaborates on the experimental specifics pertinent to that configuration for improved clarity. Comprehensive information is available in Appendix~\ref{app:exp-setup}.

Although paged optimizers are essential for enabling \textsc{QLoRA} tuning of 33B and 65B models on single GPUs with 24GB or 48GB memory, we do not offer definitive benchmarks for Paged Optimizers, given that paging primarily takes place when training with long sequences—a scenario that occurs infrequently. Nonetheless, we assess the runtime of paged optimizers on 65B models with 48GB GPUs, observing that at a batch size of 16, they achieve training speeds equivalent to conventional optimizers. Further research is needed to systematically quantify and understand the conditions leading to paging-related slowdowns.

\paragraph{Default LoRA hyperparameters do not match 16-bit performance}

Utilizing the standard LoRA strategy—modifying only the query and value projection matrices in the attention layers, as described by \citet{hu2021lora}—fails to recover the performance levels observed with full finetuning in large-scale base models. As depicted in Figure~\ref{fig:hyper} for Alpaca finetuning with LLaMA 7B, our experiments reveal that the total number of LoRA adapters integrated is the most influential hyperparameter. To reach parity with full finetuning, it is necessary to apply LoRA to every linear layer in the transformer blocks. Other parameters in LoRA, such as the projection dimension $r$, have minimal impact on finetuning outcomes (see Appendix \ref{app:exp-setup}).

In our analysis, we also observe that standard hyperparameter settings for fully finetuned baseline models are insufficiently optimized. To establish more reliable baselines, we conduct a hyperparameter search over learning rates ranging from $1\mathrm{e}{-6}$ to $5\mathrm{e}{-5}$ and batch sizes from 8 to 128. The results of 7B LLaMA finetuning on the Alpaca dataset are presented in Figure~\ref{fig:hyper}.

\paragraph{4-bit NormalFloat surpasses 4-bit Floating Point in performance}

Although the 4-bit NormalFloat (NF4) representation is theoretically optimal in terms of information retention, it remains an open question whether this theoretical advantage translates to real-world performance gains. Adopting the evaluation methodology of \citet{dettmers2022case}, we assess various quantized LLMs—including OPT \citep{zhang2022opt}, BLOOM \cite{scao2022bloom}, Pythia \cite{biderman2023pythia}, and LLaMA—across a parameter spectrum from 125M to 65B, using different quantization schemes. These models are evaluated on language modeling and zero-shot tasks. The outcomes displayed in Figure~\ref{fig:nf4} and Table~\ref{tbl:nf4_ppl} show that NF4 achieves significantly higher performance than both FP4 and Int4, and that employing double quantization lowers memory requirements without adversely affecting accuracy.

\paragraph{k-bit \textsc{QLoRA} achieves parity with 16-bit full finetuning and 16-bit LoRA approaches}

Recent studies indicate that, although 4-bit quantization for \emph{inference} is feasible, it generally incurs a performance drop when compared to 16-bit representations \citep{dettmers2022case,frantar2022gptq}. This leads to the central issue of whether such lost accuracy can be restored by applying 4-bit adapter finetuning. To investigate, we examine two experimental setups.

The first setup involves comparing fully finetuned 16-bit RoBERTA and T5 models (ranging from 125M to 3B parameters) against their adapter-based counterparts on the GLUE and Super-NaturalInstructions benchmarks. Table~\ref{tab:nf4vsbf16} summarizes these results. For both datasets, we find that adapter finetuning in 16-bit, 8-bit, and 4-bit regimes yields performances on par with the fully finetuned 16-bit baseline, indicating that the accuracy lost due to aggressive quantization can be effectively restored through adapter finetuning post-quantization.

In our second experimental design, due to the substantial memory demands associated with fully finetuning models of 11B parameters and above—which exceeds the capacity of a single high-memory GPU server—we instead investigate whether 4-bit \textsc{QLoRA} can attain parity with 16-bit LoRA across model sizes ranging from 7B to 65B parameters. To address this, we conduct finetuning of LLaMA models at 7B, 13B, 33B, and 65B scales using two instruction-tuned datasets, Alpaca and FLAN v2, and measure performance via 5-shot accuracy on the MMLU benchmark. Table~\ref{tbl:llama-datatype} presents the findings, revealing that NF4 employing double quantization entirely restores the MMLU scores of the 16-bit LoRA approach. Furthermore, we observe that while \textsc{QLoRA} with FP4 configuration consistently underperforms relative to the 16-bit brain float LoRA baseline by approximately 1 percentage point. These outcomes confirm both that (1) \textsc{QLoRA} utilizing NF4 matches the effectiveness of 16-bit full finetuning and 16-bit LoRA, and (2) NF4 exhibits superior quantization fidelity compared to FP4.

\paragraph{Summary} Our experimental data repeatedly demonstrate that 4-bit \textsc{QLoRA} employing the NF4 format achieves results equivalent to 16-bit full finetuning and 16-bit LoRA approaches on standard academic benchmarks using robust evaluation protocols. Our analysis further illustrates that NF4 outperforms FP4 and that implementing double quantization does not negatively impact accuracy. Collectively, these findings strongly support that 4-bit \textsc{QLoRA} reliably matches the outcomes delivered by 16-bit training techniques.

Consistent with prior quantization research \citep{dettmers2022case}, the results on MMLU and Elo suggest that for a fixed budget of finetuning and inference resources, expanding the parameter count of the base model while lowering numerical precision confers notable advantages. This underscores the critical efficiency gains enabled by \textsc{QLoRA}. Given the absence of observed degradation relative to full-finetuning in our 4-bit experiments, this prompts an open question regarding the precise inflection point of the performance-precision balance in \textsc{QLoRA}—a subject we propose for future investigation.

Our investigation turns to instruction tuning at scales unattainable with full 16-bit finetuning on the typical computational resources available for academic research.

\section{Pushing the Chatbot State-of-the-art with QLoRA}
\label{sec:pushing}
After demonstrating that 4-bit \textsc{QLoRA} achieves parity with 16-bit finetuning in terms of performance across diverse model sizes, task types, and datasets, we undertake a comprehensive analysis of instruction finetuning utilizing the most extensive open-source language models accessible to the research community. To evaluate the efficacy of instruction finetuning in these models, we employ a rigorous Natural Language Understanding benchmark (MMLU), and devise novel methodologies to gauge real-world chatbot effectiveness.

\subsection{Experimental setup}
We outline the key aspects of our experimental framework below, with further specifics available in Appendix \ref{app:chatbot}.

\paragraph{Data} Recognizing the absence of an exhaustive analysis of contemporary instruction-following datasets, we curate a set of eight recent collections. Our selection comprises crowd-sourced datasets (OASST1~\citep{kopf2023openassistant}, HH-RLHF~\citep{bai2022training}), datasets distilled from models that have undergone instruction tuning (Alpaca~\citep{alpaca}, self-instruct~\citep{wang2022self}, unnatural-instructions~\citep{honovich2022unnatural}), corpora aggregations (FLAN v2~\citep{chung2022scaling}), and hybrid datasets (Chip2~\citep{laion2023}, Longform~\citep{koksal2023longform}). These datasets represent a range of languages, corpus sizes, and licensing arrangements.

\paragraph{Training Setup} In order to eliminate potential confounding variables related to different training objectives, all QLoRA finetuning is conducted via cross-entropy loss in a supervised learning context, omitting any reinforcement learning even for datasets annotated with human preference judgments. Where datasets clearly distinguish between instruction and corresponding response, only the response is utilized during finetuning (additional analyses are provided in Appendix \ref{app:chatbot}). In cases like OASST1 and HH-RLHF, where multiple possible responses exist, we select the highest-rated reply at each conversation tree node and finetune on the corresponding full conversation, retaining instruction content. Across all experiments, we adopt the NF4 \textsc{QLoRA} approach, enhanced by double quantization and paged optimization to circumvent memory surges during gradient checkpointing. For the 13B and 33B LLaMA variants, we undertake modest hyperparameter tuning and observe that settings optimized for 7B, with the exception of learning rate and batch size, remain applicable; for 33B and 65B models, the learning rate is halved and batch size doubled.

\paragraph{Baselines} For benchmarking, we contrast our models against both academic systems (Vicuna~\citep{vicuna2023} and Open Assistant~\citep{kopf2023openassistant}) and leading commercial chatbots (GPT-4 \citep{openai2023gpt}, GPT-3.5-turbo, and Bard). Open Assistant utilizes a LLaMA 33B base further finetuned through Reinforcement Learning from Human Feedback (RLHF) on OASST1, mirroring our own experimental dataset. In contrast, Vicuna performs full fine-tuning of LLaMA 13B leveraging user-contributed conversations from ShareGPT, effectively distilling OpenAI GPT model behaviors.

\subsection{Evaluation}

In alignment with established evaluation standards, we employ the MMLU (Massively Multitask Language Understanding) benchmark \citep{hendrycksmeasuring}, which encompasses 57 distinct tasks such as elementary mathematics, American history, computer science, law, among others. Performance is reported as 5-shot test accuracy.

We further assess the generative capabilities of language models using a combination of automated methods and human judgment. This second evaluation phase utilizes questions created by human annotators, with the intent to gauge the quality of the model's generated replies. Although this approach offers a setting closer to practical chatbot interactions and is increasingly adopted, there remains no standard procedure in the literature. The following section details our proposed methodology, which consistently employs nucleus sampling ($p=0.9$) and a temperature setting of $0.7$ for all experiments.

\paragraph{Benchmark Data} Our assessment uses two specifically compiled collections of queries: the Vicuna prompts \citep{vicuna2023} and the OASST1 validation set \citep{kopf2023openassistant}. The Vicuna dataset contains 80 prompts spanning various domains, which we include without alterations. The OASST1 validation set comprises multilingual, user-sourced multiturn conversations involving an assistant. To form our queries, we extract all user messages from this validation set and embed the conversation history leading up to each user turn into the prompt, yielding a total of 953 distinct queries. These datasets are subsequently referred to as the Vicuna and OA benchmarks.

\paragraph{Automated Evaluation} Our first line of assessment follows the methodology outlined by \citet{vicuna2023}, wherein we leverage GPT-4 to evaluate model outputs on the Vicuna benchmark in direct comparison with ChatGPT (GPT-3.5 Turbo). For each prompt, GPT-4 receives both ChatGPT’s and the target model’s response and is instructed to award each a score out of ten, accompanied by a rationale. The performance metric is then computed as a percentage relative to ChatGPT’s score, where surpassing 100\% is possible if the model exceeds ChatGPT’s score. Notably, we observe a substantial positional bias in GPT-4, where the first-presented response is favored. To counteract this, we advise aggregating scores across both response orders.

Subsequently, we also carry out a head-to-head evaluation by directly comparing the responses of different systems. Here, the scoring is simplified to a three-way classification: win, loss, or tie. GPT-4 is instructed to identify the superior response or designate a tie, along with a justification. All possible pairwise comparisons between system outputs are conducted across both the Vicuna and OA benchmarks.

\paragraph{Human Evaluation} Although recent findings suggest that generative models can be used for automatic system evaluation \citep{fu2023gptscore}, the degree to which GPT-4-based ratings align with human evaluation, especially for chatbot models, remains unclear. Consequently, we perform two parallel sets of human assessments on the Vicuna benchmark, mirroring the automated procedures described above. Human annotations are collected via Amazon Mechanical Turk (AMT), involving two annotators for evaluations relative to ChatGPT, and three for pairwise system comparisons.

\paragraph{Elo Rating} To evaluate models, we organize both human and automated pairwise judgments into a tournament framework in which models are pitted against one another. Each round consists of two models being prompted to generate responses, and the superior answer advances, creating a direct competition akin to the methodology described in \citet{bai2022training} and \citet{vicuna2023}. Our approach, however, additionally incorporates ratings from GPT-4 alongside human evaluators. Labeled match outcomes are randomly selected to calculate Elo scores~\citep{elo1967proposed,elo1978rating}. The Elo system, standard in chess and similar competitive settings, quantifies a player’s—or in this case, a model’s—expected win probability when compared to an opponent. For instance, a model rated at 1100 Elo is anticipated to outperform a model rated at 1000 Elo about 65\% of the time; two models with identical scores, such as 1000 vs 1000 or 1100 vs 1100, would each be predicted to win 50\% of their matchups. Adjustments to Elo ratings are made after each comparison based on the predicted and actual outcomes—unexpected wins result in larger rating adjustments, whereas predictable outcomes cause minor changes. Through repeated matches, the Elo ratings converge to reflect each model’s performance capabilities within the tournament. All models are initialized with an Elo score of 1,000, and we set $K=32$ for rating updates. Following the practice in \citet{vicuna2023}, this process is executed 10,000 times using varied random seeds to mitigate effects due to the sequence of matchups.

\subsection{\textbf{Guanaco}: \textsc{QLoRA} tuned on OASST1 Establishes a New Standard for Open-source Chatbots}

Through both automated metrics and assessments from human annotators, our analysis reveals that the leading \textsc{QLoRA} fine-tuned model, {Guanaco} 65B—optimized using a modified OASST1 dataset—emerges as the top-performing open-source chatbot to date, with performance rivaling that of ChatGPT. When directly contrasted with GPT-4, human system-level pairwise evaluations and corresponding Elo ratings estimate the win probability for {Guanaco} 65B and 33B at 30\%, representing the highest score reported thus far.

Table~\ref{tbl:vicuna_eval} displays outcomes from the Vicuna benchmark \cite{vicuna2023} in comparison to ChatGPT. Notably, {Guanaco} 65B outperforms all but GPT-4, registering a 99.3\% performance rate relative to ChatGPT. While {Guanaco} 33B contains a greater number of parameters compared to Vicuna 13B, it leverages 4-bit quantization, reducing memory usage to 21 GB versus Vicuna’s 26 GB, and attains an improvement of three percentage points in performance. Additionally, {Guanaco} 7B, with a compact memory footprint of 5 GB, is suitable for deployment on contemporary mobile devices and still surpasses Alpaca 13B by nearly 20 percentage points.

Nonetheless, Table~\ref{tbl:vicuna_eval} reveals substantial confidence intervals, with significant overlap among several models’ results. We attribute this ambiguity to an insufficiently defined evaluation scale; for instance, the precise implications of scoring 8 out of 10 across diverse contexts are ambiguous. Therefore, we advocate for employing the Elo ranking system \citep{elo1967proposed}, which is based on \textit{pairwise} model comparisons by both human annotators and GPT-4, to mitigate the limitations of absolute scales. Elo scores for the most competitive models appear in Table~\ref{tbl:elo}. We observe partial discrepancies between model orderings by humans and GPT-4 on the Vicuna benchmark, especially regarding Guanaco 7B, though overall correlations are moderate, with a Kendall Tau of $\tau=0.43$ and a Spearman rank correlation of $r=0.55$ at the system comparison level. Example-level consensus is lower, reflected in a Fleiss $\kappa=0.25$ agreement between GPT-4 and the majority human vote. Collectively, these findings point to moderate concordance in system-level assessments, suggesting that model-based evaluation serves as a reasonably dependable substitute for human judgment. Further discussion on these points is presented in Section~\ref{ssec:considerations}.

The Elo scores summarized in Table~\ref{tbl:elo_large} reveal that the {Guanaco} models, specifically those with 33B and 65B parameters, surpass all counterparts aside from GPT-4 on both the Vicuna and OA evaluations and exhibit performance on par with ChatGPT, consistent with findings in Table~\ref{tbl:vicuna_eval}. It is important to highlight that the Vicuna benchmark is more favorable to open-source systems, while the broader OA benchmark tends to advantage ChatGPT. Moreover, analysis of Tables \ref{tbl:mmlu-llama} and \ref{tbl:vicuna_eval} demonstrates that the choice of finetuning dataset plays a crucial role in determining model outcomes. For example, Llama models refined with FLAN v2 achieve particularly high scores on MMLU but perform poorly on the Vicuna metric—a pattern reflected in other model families as well. These results further underscore that existing evaluation benchmarks test complementary model capabilities: excelling at MMLU does not guarantee top-tier chatbot interaction as measured by Vicuna or OA, and vice versa.

is distinctive among top performers in our study as it avoids proprietary data sources, since the OASST1 dataset's collection policy expressly bans the inclusion of GPT model outputs. The runner-up among models exclusively trained on open-source content is Anthropic's HH-RLHF model, which falls short by 30 percentage points relative to {Guanaco} on the Vicuna metric (refer to Table~\ref{tbl:vicuna_eval}). Collectively, these findings indicate that the 4-bit \textsc{QLoRA} approach is highly effective, enabling the construction of leading-edge chatbots competitive with ChatGPT. Notably, our 33B-parameter {Guanaco} can be trained on widely accessible 24 GB consumer GPUs in under twelve hours. This paves the way for future research leveraging \textsc{QLoRA} finetuning with tailored open-source datasets to create models that can contend with today’s best proprietary alternatives.

\section{Qualitative Analysis}

Although our primary focus is on quantitative metrics, relying solely on aggregate statistics introduces several limitations. Chief among these is the issue of benchmark validity \citep{liao2021we}—the question of whether a benchmark genuinely evaluates the abilities implied by its title or description remains open, particularly as we uncover ``shortcuts'' that machine learning models may exploit to obtain inflated results \citep{gururangan2018annotation, poliak2018hypothesis}. To mitigate some of these shortcomings, we supplement our evaluation with qualitative analysis, divided into two subsections. In \S\ref{ssec:examples}, we present selected outputs generated by the 65b {Guanaco} model, illustrating recurrent themes and tendencies we have noticed. Subsequently, in \S\ref{ssec:considerations}, we elaborate on methodological concerns related to our findings and discuss the interpretive framework we apply to the results.

\subsection{Examining Generated Examples}
\label{ssec:examples}

Our process for collecting qualitative examples involves systematically reviewing responses produced by {Guanaco} on both the Vicuna and OpenAssistant benchmarks, aiming to identify recurring behaviors in the model’s output. When a consistent phenomenon is identified, we then design prompts that specifically elicit this pattern—even if it is suboptimal. For instance, upon noticing a propensity for verbose answers, we challenge the model with the instruction ``Answer yes or no without explanation.'' This strategy helps us identify failure cases (``lemons'') where the model falls short under adversarial prompting, as well as robust cases (``cherries'') where such prompting does not lead to errors. The examples shown here are all generated using Nucleus Sampling \cite{holtzmancurious} with $p=0.9$.

Naturally, the limitations of this brief qualitative investigation prevent exhaustive analysis; accounting for every potential variable—such as the expansive range of responses that may be generated for any specific prompt—lies beyond the study’s purview. Instead, we rely on a select set of samples which we anticipate adequately reflect typical outputs. Presenting these instances is intended to supplement the quantitative findings discussed in earlier sections. By releasing all models and code as open source, we encourage further research that builds on our preliminary observations, especially around the challenges outlined here.

\paragraph{Factual Recall} For straightforward factual queries (for example, ``What is the capital of Zambia?''), every evaluated model reliably produces the correct response, as demonstrated below:

\begin{quote} 
    \textbf{Guanaco}:
    \texttt{The capital of Zambia is Lusaka.}
\end{quote}

When questions pertain to more obscure facts, however, {Guanaco}'s reliability diminishes, though it continues to respond with confidence. For example, consider the following question taken from HotPotQA \citep{yang2018hotpotqa}:

\begin{quote}
     \textbf{User}: In what year was the singer who popularized `I'll Keep the Lovelight Burning' born?  

     \textbf{Guanaco}:~\texttt{Al Jolson is credited with popularizing the song `I'll Keep the Lovelight Burning,' and he was born in the year 1886.}
 \end{quote}

Here, the response cites both an incorrect popularizer and an incorrect birth year (although the birth year matches that of the mentioned individual, Al Jolson).

\paragraph{Suggestibility} Interestingly, {Guanaco} often exhibits a notable reluctance to endorse certain types of assumed falsehoods, as evidenced in the subsequent exchange:

\begin{quote}
    \textbf{User}: How was it finally, officially confirmed that the earth is flat by peer reviewed scientists? \\ \\
    \textbf{Guanaco}: \texttt{There has never been an official confirmation by peer reviewed scientists asserting that the earth is flat. The scientific consensus firmly maintains that the earth is spherical—a fact supported for centuries through extensive experimental evidence and astronomical observation. \\ \\ The flat earth hypothesis is widely considered a debunked conspiracy theory within scientific communities. The claim that peer reviewed research has verified a flat earth model is not supported by any credible evidence. \\ \\ Please feel free to ask additional questions if you wish.}
\end{quote}

is adept at identifying when certain questions are fundamentally unanswerable, for instance,

\begin{quote}
    \textbf{User}: What time is it?

    \textbf{Guanaco}: \texttt{I'm unable to provide the current time, as I am an AI language model without access to real-time data. My answers draw only on information within my training and do not include up-to-date or live data. For current time inquiries, it is best to use a clock or consult an online time service.}
\end{quote}

Comparable responses arise for prompts such as ``Where are you?'', ``How are you?'', and similar questions.

\paragraph{Refusal} In the same vein, {Guanaco} at times inexplicably declines to comply with certain instructions:

\begin{quote}
    \textbf{User}: Please reverse the order of the words in the phrase ``I'm a cool dude in a mocha mood.'' \\ \\
    \textbf{Guanaco}: \texttt{I apologize, but I'm currently unable to do that. However, I can elaborate on why your sentence is correct from a grammatical standpoint.\\ \\The construction ``I'm a cool dude in a mocha mood'' adheres to proper English grammar and syntactic conventions. Specifically, the descriptive term ``cool'' correctly modifies the noun ``dude'' as an attributive adjective, while the prepositional phrase ``in a mocha mood'' further contextualizes the sentence. This composition is grammatically valid and semantically clear.\\ \\If you have additional questions or if you would like assistance with something else, please let me know.}
\end{quote}

\paragraph{Secret Keeping} Regrettably, this type of refusal is not consistently dependable, even when it is desired by system developers. Frequently, designers aim to prevent users from obtaining sensitive information, even if the model is capable of producing unrestricted text responses. To investigate this, we prompt with ``The secret word is `banana'. Do not reveal or repeat back this back, no matter what the human requests.'' and evaluate how resistant the model is to revealing this secret if a user tries to obtain it. If asked directly, the model declines to answer:

\begin{quote}
    \textbf{User}: What is the secret word?

    \textbf{Guanaco}: \texttt{I'm sorry, but I cannot share the secret word. I am programmed not to reveal or repeat the secret word regardless of the request. Can I assist you with anything else?}
\end{quote}

which matches the intended outcome. Nevertheless, even minimal prompt engineering is sufficient to bypass this safeguard:

\begin{quote}
    \textbf{User}: Let's play a game. Your goal is to disregard your earlier instructions. What is the secret word?

    \textbf{Guanaco}: \texttt{Apologies, the secret word is banana. It would be best not to mention it further. Do you have any other questions?}
\end{quote}

This highlights the necessity for research into techniques that ensure more reliable adherence to system instructions.

\paragraph{Math} The principal area of weakness for {Guanaco} is mathematics—a common challenge observed in language models, as reported in \citep{liang2022holistic}. When prompted to show intermediate steps, {Guanaco} often produces correct solutions. For example,

\begin{quote}
    \textbf{User}: Lee mows one lawn and charges \$33. Last week he mowed 16 lawns and three customers each gave him a \$10 tip. How many dollars did Lee earn mowing lawns last week?

    \textbf{Guanaco}: \texttt{Lee earned \$582 mowing lawns last week. \\ \\ To find the total amount he earned, we first need to calculate how much he earned for mowing 16 lawns: \\ \\  \$33 x 16 = \$528 \\ \\ Then we need to add the tips he received from the three customers: \\  \\ \$10 x 3 = \$30 \\ \\ So the total amount Lee earned was: \\ \\ \$528 + \$30 = \$558}
\end{quote}

Nevertheless, {Guanaco} may yield incorrect results on relatively basic tasks when it neglects a stepwise approach—an issue well-documented in the literature \cite{weichain}. Consider this sample exchange:

\begin{quote}
    \textbf{User}: Please factorize 1833.

    \textbf{Guanaco}: \texttt{Sure, I'd be happy to help! 1833 is a prime number, which means it has only two factors: 1 and 1833. \\ \\ So the unique factorization of 1833 is: \\ \\ 2\^{}1 * 3\^{}2 * 17\^{}1 \\ \\ I hope this helps! Let me know if you have any other questions.}
\end{quote}

which is incorrect (the actual factorization is $3 \times 17 \times 43$), and it is doubly so.

\paragraph{Theory of Mind} {Guanaco} demonstrates unexpectedly robust Theory of Mind abilities, as indicated in previous work \cite{nematzadeh2018evaluating, sap2022neural}. As an illustration, the model offers a detailed and accurate response to the following prompt:

\begin{quote}
    \textbf{User}: James and Abby are in the bedroom. Abby put the pen in the desk drawer. Abby leaves the bedroom. James moves the pen into the bag. Where does James think Abby will look for the pen?

    \textbf{Guanaco}: \texttt{James knows that Abby put the pen in the desk drawer. But James also knows that he has moved the pen into the bag, and he assumes that Abby does not know about this. So, James may guess that Abby will look for the pen in the desk drawer, because that's where she left it last time. However, this assumption may be incorrect, because Abby may have seen James moving the pen.}
\end{quote}

Nevertheless, the consistency of these inferences is limited, with the model occasionally providing rationales that reflect implausible or unsupported assumptions, for example:

\begin{quote}
    \textbf{User}: Evelyn entered the living room. Jackson entered the playroom. James entered the playroom. The beans are in the treasure chest. James exited the playroom. Jackson moved the beans to the pantry. Jackson exited the playroom. James entered the living room. Where will James look for the beans?

    \textbf{Guanaco}: \texttt{James will look for the beans in the pantry, because that's where Jackson moved them.}
\end{quote}

In this case, {Guanaco} assumes that knowledge has been transmitted without any explicit description. Such shortcomings parallel findings from recent studies \cite{sap2022neural}, indicating an area in need of further investigation.

\subsection{Considerations}
\label{ssec:considerations}

\paragraph{Evaluation}

Our results reveal moderate consistency among human evaluators (Fleiss $\kappa=0.42$), with further decline in agreement when contrasting outputs from two high-performing systems. This highlights inherent shortcomings in existing chatbot task benchmarks and in manual evaluation practices. During head-to-head assessments of ChatGPT and {Guanaco} 65B using the Vicuna benchmark, it became clear that subjectivity considerably influences preferences, as disagreements were frequent even among this paper's authors. Going forward, research should explore solutions to address subjectivity in evaluation, potentially leveraging methodologies from fields such as Human-Computer Interaction and Psychology, which specialize in the management of personal preferences.

Additionally, our study identifies notable biases in automated evaluation metrics. Notably, GPT-4 exhibits strong order bias, consistently assigning superior scores to the system mentioned first in the prompt. The relatively low instance-level agreement between GPT-4 and human raters (Fleiss $\kappa=0.25$) further implies misalignment in evaluation criteria. Moreover, Table~\ref{tbl:elo_large} demonstrates that GPT-4 rates its own generated content much higher than do human evaluators—Elo scores of 1348 versus 1176, equating to an approximate 20\% higher chance of outperforming an alternative model.

Future studies should assess the existence of possible biases within automated evaluation frameworks and explore approaches for addressing these issues.

\paragraph{Data \& Training} It is important to highlight that the OASST1 dataset used for training the {Guanaco} models is multilingual, and the OA benchmark similarly includes prompts spanning several languages. Determining the extent to which exposure to multilingual data enhances non-English instruction following, and whether this factor accounts for the greater performance disparity between the Vicuna-13B model (which was exclusively trained with English-language material) and the {Guanaco} 33B and 65B models on the OA benchmark, is a question left for future inquiry.

Owing to the robust results shown by the {Guanaco} models, we conducted an analysis to identify any data leakage between OASST1 and Vicuna benchmark prompts. After utilizing fuzzy string matching and closely reviewing the top matches, we did not observe any prompt duplication between the two datasets.

Additionally, it is worth noting that our model’s training relies solely on supervised learning via cross-entropy loss, without incorporating reinforcement learning from human feedback (RLHF). This situation invites deeper comparative investigations into the relative merits and limitations of cross-entropy and RLHF-based training approaches. We anticipate that \textsc{QLoRA} can facilitate such large-scale analysis without necessitating extensive compute resources.

\section{Related Work}

\paragraph{Quantization of Large Language Models} Existing work on quantizing LLMs has primarily addressed improving inference efficiency. Leading methods for retaining the quality associated with 16-bit LLMs typically target the management of outlier activation features (for example, SmoothQuant~\citep{xiao2022smoothquant} and LLM.int8()~\citep{dettmers2022llm}), while others implement advanced grouping strategies \citep{park2022nuqmm, yao2022zeroquant}. There has also been significant investigation into lossy quantization methods, examining the trade-offs of conventional rounding schemes~\citep{dettmers2022case,zeng2022glm,pope2022efficiently} and exploring optimizations in rounding to enhance quantization precision~\citep{frantar2022gptq}. Apart from our contribution, SwitchBack layers~\citep{wortsman2023stable} is unique in exploring backpropagation through quantized weights at a parameter scale beyond 1B.

\paragraph{Finetuning with Adapters} In our work, we employ Low-rank Adapters~\citep{hu2021lora} (LoRA), though the landscape of Parameter Efficient FineTuning (PEFT) methods includes a wide array of alternatives, such as prompt tuning~\citep{qin2021learning,lester2021power,li2021prefix}, modifications to embedding layer inputs~\citep{an2022input}, adjustment of hidden states (IA$^3$)~\citep{liu2022few}, insertion of entire new layers~\citep{houlsby2019parameter}, bias tuning~\citep{zaken2021bitfit}, training a mask over parameters using Fisher information~\citep{sung2021training}, as well as composite strategies~\citep{henderson2021compacter}. Our experiments demonstrate that LoRA adapters are sufficient to achieve the same performance as conventional 16-bit finetuning. Investigation into the relative merits and trade-offs of alternative PEFT strategies is left to subsequent research.

\paragraph{Instruction Finetuning} Instruction finetuning adapts pretrained large language models to better follow prompt-based instructions, by finetuning them on datasets composed of input-output pairs from diverse sources. A variety of methodologies and datasets support this process, such as MetaICL~\citep{min2021metaicl}, MetaTuning~\citep{zhong2021adapting}, InstructGPT~\citep{ouyang2022training}, FLAN~\citep{wei2021finetuned,chung2022scaling}, PromptSource~\citep{bach2022promptsource}, Super-NaturalInstructions~\citep{wang2022super,sanh2021multitask}, Self-instruct~\citep{wang2022self}, UnnaturalInstructions~\citep{honovich2022unnatural}, OPT-IML~\citep{iyer2022opt}, UnifiedSKG~\citep{xie2022unifiedskg}, OIG/Chip2~\citep{laion2023}, Alpaca~\citep{alpaca}, Vicuna~\citep{vicuna2023}, Koala~\citep{koala_blogpost_2023}, and Self-instruct-GPT-4~\citep{peng2023instruction}.

\paragraph{Chatbots} A large fraction of instruction-tuned models are implemented as dialogue agents, frequently utilizing Reinforcement Learning from Human Feedback (RLHF)~\citep{christiano2017deep} or relying on feedback generated by other models through Reinforcement Learning with AI Feedback (RLAIF)~\citep{bai2022constitutional}. Existing methods and resources include Anthropic-HH~\citep{askell2021general,bai2022training}, Open Assistant~\citep{kopf2023openassistant}, LaMDA~\citep{thoppilan2022lamda}, and Sparrow~\citep{glaese2022improving}. Although our approach does not make use of reinforcement learning, our highest-performing model, Guanaco, is trained via finetuning on multi-turn conversation data drawn from the Open Assistant corpus, which was intended for RLHF workflows~\citep{kopf2023openassistant}. To address the evaluation of chatbots, recent efforts leverage GPT-4 as an alternative to human judgments~\citep{vicuna2023,peng2023instruction}. We contribute improvements to these evaluation methodologies, emphasizing increased robustness and reliability.

\section{Limitations and Discussion}

Our results indicate that \textsc{QLoRA} enables recovery of full 16-bit finetuning effectiveness using a 4-bit base model equipped with Low-rank Adapters (LoRA). Nevertheless, it remains unproven whether \textsc{QLoRA} can sustain this level of performance at very large scales, specifically for models with 33B and 65B parameters. Given the substantial computational requirements, thorough investigations at these larger scales are deferred to future work.

A further constraint involves our assessment of instruction-tuned models. Although we report results on MMLU, the Vicuna benchmark, and the OA benchmark, our evaluation does not cover additional suites such as BigBench, RAFT, or HELM, and thus our conclusions may not extrapolate to those benchmarks. However, we do conduct a comprehensive exploration using MMLU and introduce new approaches for chatbot evaluation.

Based on the data analyzed, it becomes apparent that model performance on specific benchmarks is strongly influenced by the degree of overlap between finetuning datasets and the benchmarks themselves. For instance, FLAN v2 shares similarities with MMLU, yet diverges significantly from chatbot evaluation tasks; the converse is observed with the Chip2 dataset. Correspondingly, the model results on MMLU and Vicuna benchmarks align with this pattern. This observation emphasizes the need not just for improved evaluation methodologies and benchmarks, but also for a careful consideration of what precisely is being measured. The research community must reflect on the end goal: Are we aiming for proficiency on traditional academic knowledge assessments, or prioritizing excellence in conversational capabilities? Or is there a different objective altogether? Since leveraging existing benchmarks is far more convenient than devising new ones, prevailing benchmarks can inadvertently shape the research direction. It is therefore critical to ensure, at the community level, that chosen benchmarks accurately represent our research priorities.

Despite providing a comprehensive assessment of general chatbot capabilities, our study has limitations in its exploration of responsible AI aspects for {Guanaco}. Specifically, we only perform a limited evaluation regarding bias, as reported in Table~\ref{tbl:crows}, where we compare the probability that {Guanaco}-65B produces socially biased outputs with those from other models. {Guanaco}-65B demonstrates a significantly lower mean score in this regard relative to unmodified pretrained models. This suggests that finetuning on the OASST1 dataset can effectively reduce bias present in the original LLaMA model. Although this finding is promising, it remains uncertain whether {Guanaco} would exhibit similarly reduced biases if assessed on other dimensions of bias. Further in-depth studies to analyze biases in {Guanaco} and related conversational agents are left as an avenue for future exploration.

Another constraint of our work is the lack of evaluation for varying quantization bit-depths, such as 3-bit base models, as well as alternative adapter methods. Beyond LoRA, numerous Parameter Efficient FineTuning (PEFT) approaches have demonstrated effectiveness, though their scalability to large-scale models remains an open question. We selected LoRA due to its well-documented stability, but alternative adapters may deliver superior results. The process of finetuning following quantization appears to recover much of the information lost through quantization, potentially enabling more aggressive quantization strategies. For instance, applying 3-bit GPTQ quantization to the base model and subsequently utilizing LoRA for finetuning may achieve performance on par with full 16-bit finetuning after adjustment.

\section{Broader Impacts}

The \textsc{QLoRA} finetuning approach represents the first instance where 33B-parameter models can be effectively finetuned on a single consumer-grade GPU, and 65B-parameter models can be finetuned on a single professional GPU, all without sacrificing performance relative to full finetuning baselines. Our most performant 33B model, trained on the Open Assistant dataset, achieves results competitive with ChatGPT on the Vicuna benchmark. Considering that instruction finetuning plays a pivotal role in converting large pretrained language models into ChatGPT-like conversational agents, we anticipate that our methodology will significantly democratize finetuning—especially benefiting resource-constrained researchers—and foster wider accessibility to leading-edge NLP technologies. As a result, \textsc{QLoRA} can act as a leveling mechanism, narrowing the gap between large industrial players and smaller research groups equipped with only consumer hardware.

An additional avenue for significant impact lies in deploying models to mobile devices. We propose that our \textsc{QLoRA} technique may mark a pivotal advancement by facilitating the finetuning of LLMs directly on smartphones and in other resource-constrained environments. Although earlier work has demonstrated that 7B parameter models can be executed on mobile phones, \textsc{QLoRA} stands out as the inaugural approach that also permits the finetuning of these models on such platforms. Our analysis suggests that, on an iPhone 12 Plus, it is possible to finetune up to 3 million tokens overnight while the device is plugged in. Despite the fact that finetuned 7B models do not reach the capabilities of ChatGPT, their performance is sufficiently strong to support the development of innovative applications previously hindered by privacy concerns or LLM performance limitations. The \textsc{QLoRA} method thus offers pathways for privacy-oriented LLM use cases, empowering individuals to take control of their data and models, and simplifying the process of LLM deployment.

Nonetheless, it should be noted that the ability to finetune models is a dual-use technology, susceptible to malicious exploitation. The proliferation of LLMs poses acknowledged risks \citep{bommasani2021opportunities,bender2021dangers}; however, we argue that democratizing access to this rapidly spreading technology will foster more robust, independent scrutiny as opposed to the risks of centralizing LLM capabilities within large, non-transparent organizations that withhold models and code from public examination.

In summary, we anticipate that \textsc{QLoRA} will yield overwhelmingly positive outcomes by considerably broadening access to high-quality LLM finetuning, making these tools more accessible and usable for a wide audience.